% Section 5 targets. This file is included by main.tex.

\setcounter{section}{4}
\section{QUBO Transformation}

Section 4 of the paper is an informal overview, so it introduces no Lea target.
This section formalizes the transformation itself. Its objective is first
defined directly and then represented by a QUBO weight matrix. This separation
makes the correctness proof independent of coefficient-collection details.

\subsection{Capacity encoding}

The paper requires storage bits whose selected weights can represent every
amount from zero through a node's capacity. The following abstraction records
exactly that property.

\begin{definition}[Capacity encoding]\label{def:capacity-encoding}
% lea: define
% lea: label=CapacityEncoding
% lea: context={Create a structure CapacityEncoding (C : Nat). It should contain a finite linearly ordered chunk type, a weight : Chunk -> Nat, a proof that the sum of all weights is C, an encode function taking t : Nat and h : t <= C to Chunk -> Fin 2, and a proof that the weighted value of encode t h equals t. Also define a weightedValue helper.}
For $C\in\mathbb{N}$, a \emph{capacity encoding} consists of a finite linearly
ordered chunk type $J_C$, weights $\omega_C:J_C\to\mathbb{N}$ satisfying
\[
  \sum_{i\in J_C}\omega_C(i)=C,
\]
and, for every $t\leq C$, a binary vector
$\operatorname{encode}_C(t):J_C\to\{0,1\}$ satisfying
\[
  \sum_{i\in J_C}\omega_C(i)\operatorname{encode}_C(t)(i)=t.
\]
Because the weights are natural numbers and sum to $C$, every binary vector
over $J_C$ represents a value at most $C$.
\end{definition}

\begin{definition}[Corrected binary capacity encoding]
  \label{def:binary-capacity-encoding}
% lea: define
% lea: label=binaryCapacityEncoding
% lea: uses={CapacityEncoding}
% lea: context={Define binaryCapacityEncoding (C : Nat) : CapacityEncoding C. For C = 0 use no chunks. For C > 0 let m = Nat.log2 C, use weights 2^j for j < m, and use final weight C - (2^m - 1). Prove the weights sum to C and every t <= C is representable. This gives Nat.log2 C + 1 chunks when C > 0 and reproduces weights 1,2,4 for C=7 and 1,2 for C=3.}
For $C=0$, use no chunks. For $C>0$, let
$m=\lfloor\log_2 C\rfloor$ and use one weight $2^j$ for every $0\leq j<m$,
together with the final weight
\[
  C-(2^m-1).
\]
The lower powers represent every value from $0$ through $2^m-1$. The final
weight is at most $2^m$, so adding it produces the remaining values without a
gap. The weights sum to $C$, giving a valid capacity encoding with $m+1$
chunks. In particular, capacities $7$ and $3$ yield $(1,2,4)$ and $(1,2)$.
\end{definition}

\subsection{Variables and decoding}

\begin{definition}[Allocation QUBO variables]\label{def:allocation-variables}
% lea: define
% lea: label=AllocationQUBOVariable
% lea: uses={AllocationInstance, binaryCapacityEncoding}
% lea: context={Define a finite linearly ordered type AllocationQUBOVariable (I : AllocationInstance) with two disjoint constructors: assignment p n for p : I.Partition and n : I.Node, and storage n i where i is a chunk of binaryCapacityEncoding (I.capacity n). Provide the Fintype, DecidableEq, and LinearOrder instances needed by QUBO.}
For an allocation instance $I$, the variable type $V_I$ is the disjoint union
of assignment variables and storage variables:
\[
  V_I=(P_I\times N_I)\ \sqcup\
  \{(n,i)\mid n\in N_I,\ i\in J_{\operatorname{capacity}_I(n)}\}.
\]
We write $A_{pn}$ for the assignment variable associated with $(p,n)$ and
$S_{in}$ for the storage variable associated with chunk $i$ of node $n$.
\end{definition}

\begin{definition}[Decoding an assignment]\label{def:decode-assignment}
% lea: define
% lea: label=DecodeAssignment
% lea: uses={AllocationInstance, DataAssignment, AllocationQUBOVariable}
% lea: context={Define AssignmentBit and StorageBit accessors for a configuration x : AllocationQUBOVariable I -> Fin 2. Define DecodeAssignment I x : DataAssignment I by putting n in the Finset for p exactly when AssignmentBit x p n is 1. Keep all three declarations in this target artifact, with DecodeAssignment as the primary declaration.}
For a binary configuration $x:V_I\to\{0,1\}$, write $A_x(p,n)$ and $S_x(n,i)$
for the corresponding bit values. Decode $x$ into the data assignment
\[
  D_x(p)=\{n\in N_I\mid A_x(p,n)=1\}.
\]
Storage bits do not affect the decoded data assignment; they witness storage
consumption in the penalty formula.
\end{definition}

\subsection{Objective components}

All objective components are real-valued. Differences are formed in
$\mathbb{Z}$ or $\mathbb{R}$ rather than in $\mathbb{N}$, so subtraction does
not truncate.

\begin{definition}[Safety penalty]\label{def:safety-penalty}
% lea: define
% lea: label=SafetyPenalty
% lea: uses={AllocationInstance, AllocationQUBOVariable, DecodeAssignment}
% lea: context={Define SafetyPenalty (I : AllocationInstance) (x : AllocationQUBOVariable I -> Fin 2) : Real using the paper's formula. Cast the assignment-bit sum and I.k to Int or Real before subtraction, then square and sum over partitions. Prove or expose nonnegativity for later use.}
The safety penalty is
\[
  Q_R(x)=\sum_{p\in P_I}
    \left(\sum_{n\in N_I} A_x(p,n)-k_I\right)^2.
\]
It is zero exactly when every partition is assigned to exactly $k_I$ nodes.
\end{definition}

\begin{definition}[Storage penalty]\label{def:storage-penalty}
% lea: define
% lea: label=StoragePenalty
% lea: uses={AllocationInstance, AllocationQUBOVariable, DecodeAssignment, binaryCapacityEncoding}
% lea: context={Define StoragePenalty (I : AllocationInstance) (x : AllocationQUBOVariable I -> Fin 2) : Real. For each node square the difference between assigned partition size and the weighted value of its selected storage chunks, then sum over nodes. Cast before subtraction and expose nonnegativity.}
The storage penalty is
\[
  Q_S(x)=\sum_{n\in N_I}
  \left(
    \sum_{p\in P_I} A_x(p,n)\operatorname{size}_I(p)
    -\sum_{i\in J_{\operatorname{capacity}_I(n)}}
       \omega_{I,n}(i)S_x(n,i)
  \right)^2.
\]
If this penalty is zero, assigned storage equals a representable value and is
therefore at most the node's capacity. Conversely, every storage-feasible
assignment admits storage bits that make the penalty zero.
\end{definition}

\begin{definition}[Processing objective]\label{def:processing-objective}
% lea: define
% lea: label=ProcessingObjective
% lea: uses={AllocationInstance, AllocationQUBOVariable, DecodeAssignment}
% lea: context={Define ProcessingObjective (I : AllocationInstance) (x : AllocationQUBOVariable I -> Fin 2) : Real as the finite double sum of requests times communicationCost times (1 - AssignmentBit). Cast the natural coefficients and bit values to Real.}
The processing objective is
\[
  Q_C(x)=\sum_{p\in P_I}\sum_{n\in N_I}
    \operatorname{communicationCost}_I(p)
    \operatorname{requests}_I(p,n)
    \bigl(1-A_x(p,n)\bigr).
\]
\end{definition}

\begin{definition}[Remote-cost bound]\label{def:remote-cost-bound}
% lea: define
% lea: label=RemoteCostBound
% lea: uses={AllocationInstance}
% lea: context={Define RemoteCostBound (I : AllocationInstance) : Real as the sum over every partition and node of I.communicationCost p times I.requests p n, cast to Real. Prove or expose that every ProcessingObjective lies between zero and this bound.}
The cost of serving every request remotely is
\[
  B_I=\sum_{p\in P_I}\sum_{n\in N_I}
    \operatorname{communicationCost}_I(p)
    \operatorname{requests}_I(p,n).
\]
For every configuration $x$, $0\leq Q_C(x)\leq B_I$.
\end{definition}

\begin{definition}[Valid penalty weight]\label{def:valid-penalty-weight}
% lea: define
% lea: label=ValidPenaltyWeight
% lea: uses={RemoteCostBound}
% lea: context={Define ValidPenaltyWeight (I : AllocationInstance) (h : Real) : Prop to mean RemoteCostBound I < h.}
A real penalty weight $h$ is \emph{valid} for $I$ when
\[
  B_I<h.
\]
\end{definition}

\begin{definition}[Penalized allocation objective]\label{def:allocation-objective}
% lea: define
% lea: label=AllocationObjective
% lea: uses={SafetyPenalty, StoragePenalty, ProcessingObjective}
% lea: context={Define AllocationObjective (I : AllocationInstance) (h : Real) (x : AllocationQUBOVariable I -> Fin 2) : Real by h * (SafetyPenalty I x + StoragePenalty I x) + ProcessingObjective I x.}
The complete penalized objective is
\[
  Q_{I,h}(x)=h\bigl(Q_R(x)+Q_S(x)\bigr)+Q_C(x).
\]
\end{definition}

\subsection{Representation as QUBO}

Expanding the squared linear expressions produces only constant, linear, and
quadratic terms. Linear terms are represented by diagonal QUBO weights because
$x(i)^2=x(i)$ for binary variables. The constant terms do not affect minimizers,
but we track them explicitly to obtain an exact energy identity.

\begin{definition}[Objective offset]\label{def:objective-offset}
% lea: define
% lea: label=AllocationObjectiveOffset
% lea: uses={AllocationInstance, RemoteCostBound}
% lea: context={Define AllocationObjectiveOffset (I : AllocationInstance) (h : Real) : Real as h * Fintype.card I.Partition * (I.k)^2 + RemoteCostBound I, with all natural values cast to Real. This is the constant term of the expanded AllocationObjective.}
The constant part of the expanded objective is
\[
  K_{I,h}=h\,|P_I|\,k_I^2+B_I.
\]
\end{definition}

\begin{theorem}[The objective is representable as a QUBO]
  \label{thm:objective-is-qubo}
% lea: formalize
% lea: label=allocationObjective_is_qubo
% lea: uses={QUBO, AllocationQUBOVariable, AllocationObjective, AllocationObjectiveOffset}
% lea: context={Prove that for every I and h there exists q : QUBO (AllocationQUBOVariable I) such that for every binary configuration x, q.energy x + AllocationObjectiveOffset I h = AllocationObjective I h x. Construct q by expanding the squares, using x^2 = x for Fin 2 values, and collecting coefficients under the variable LinearOrder.}
For every allocation instance $I$ and real $h$, there exists a QUBO $q_{I,h}$
on $V_I$ such that every binary configuration $x$ satisfies
\[
  E_{q_{I,h}}(x)+K_{I,h}=Q_{I,h}(x).
\]
\end{theorem}

\begin{definition}[The allocation QUBO]\label{def:allocation-qubo}
% lea: define
% lea: label=AllocationQUBO
% lea: uses={allocationObjective_is_qubo}
% lea: context={Define AllocationQUBO (I : AllocationInstance) (h : Real) : QUBO (AllocationQUBOVariable I) by choosing the witness from allocationObjective_is_qubo. Also prove a named theorem allocationQUBO_energy_add_offset giving the witness identity for all x.}
Fix one QUBO supplied by Theorem~\ref{thm:objective-is-qubo} and denote it by
$\operatorname{AllocationQUBO}(I,h)$. Its energy satisfies
\[
  E_{\operatorname{AllocationQUBO}(I,h)}(x)+K_{I,h}=Q_{I,h}(x)
\]
for every configuration $x$.
\end{definition}

\subsection{Correctness of the transformation}

\begin{lemma}[Zero penalty implies feasibility]
  \label{lem:zero-penalty-feasible}
% lea: formalize
% lea: label=zeroPenaltyImpliesFeasible
% lea: uses={SafetyPenalty, StoragePenalty, DecodeAssignment, FeasibleAssignment, binaryCapacityEncoding}
% lea: context={Prove that if SafetyPenalty I x = 0 and StoragePenalty I x = 0, then FeasibleAssignment I (DecodeAssignment I x). Use the exact-cardinality consequence of Q_R = 0 and the fact that every selected chunk sum is at most capacity.}
If $Q_R(x)=0$ and $Q_S(x)=0$, then the decoded assignment $D_x$ is feasible.
\end{lemma}

\begin{lemma}[Feasible assignments have zero-penalty encodings]
  \label{lem:feasible-zero-penalty}
% lea: formalize
% lea: label=feasibleAssignmentHasZeroPenaltyEncoding
% lea: uses={FeasibleAssignment, DecodeAssignment, SafetyPenalty, StoragePenalty, binaryCapacityEncoding}
% lea: context={Prove that every feasible D : DataAssignment I has a binary configuration x with DecodeAssignment I x = D, SafetyPenalty I x = 0, and StoragePenalty I x = 0. Set assignment bits from D and obtain each node's storage bits from binaryCapacityEncoding.encode using the storage-feasibility proof.}
For every feasible data assignment $D$, there exists a binary configuration $x$
such that
\[
  D_x=D,\qquad Q_R(x)=0,\qquad Q_S(x)=0.
\]
\end{lemma}

\begin{lemma}[The processing objectives agree]
  \label{lem:processing-objectives-agree}
% lea: formalize
% lea: label=processingObjective_eq_processingCost
% lea: uses={ProcessingObjective, DecodeAssignment, ProcessingCost}
% lea: context={Prove that ProcessingObjective I x equals the real cast of ProcessingCost I (DecodeAssignment I x). Unfold the decoded Finset membership condition and the two finite sums.}
For every binary configuration $x$,
\[
  Q_C(x)=\operatorname{ProcessingCost}_I(D_x),
\]
where the natural-valued processing cost is viewed as a real number.
\end{lemma}

\begin{theorem}[Correctness of the QUBO transformation]
  \label{thm:qubo-encoding-correct}
% lea: formalize
% lea: label=quboEncodingCorrect
% lea: uses={AllocationQUBO, ValidPenaltyWeight, zeroPenaltyImpliesFeasible, feasibleAssignmentHasZeroPenaltyEncoding, processingObjective_eq_processingCost, IsOptimalAssignment}
% lea: context={Assume ValidPenaltyWeight I h and that at least one FeasibleAssignment exists. Prove both directions: every QUBO.IsOptimal (AllocationQUBO I h) x decodes to IsOptimalAssignment I (DecodeAssignment I x); and every IsOptimalAssignment I D has some encoding x that decodes to D and is QUBO.IsOptimal. Use the strict bound h > RemoteCostBound to exclude every positive integer constraint penalty, and use allocationQUBO_energy_add_offset to transfer minimizers between QUBO energy and AllocationObjective.}
Let $h$ be a valid penalty weight and suppose at least one feasible assignment
exists. Then:
\begin{enumerate}
  \item every global minimizer $x$ of $\operatorname{AllocationQUBO}(I,h)$
    decodes to an optimal feasible data assignment $D_x$;
  \item every optimal feasible data assignment $D$ has a binary encoding $x$
    with $D_x=D$ that globally minimizes
    $\operatorname{AllocationQUBO}(I,h)$.
\end{enumerate}
Thus the global minimizers of the QUBO correspond exactly, after decoding, to
the optimal solutions of the data-assignment problem.
\end{theorem}

\section*{Next formalization milestone}

After all targets above are checked, the next include file should cover the
formal analysis from Section 6: the knapsack reduction, the number of logical
variables, and dense- and sparse-topology bounds. The storage-connection count
must be audited before stating its asymptotic theorem, because the paper's
displayed bound omits the factor contributed by ranging over all nodes.
